% 기어식 구동부의 축간 거리. 단위 mm. 그림 fig:hw-gear-dim-drive
% 수치의 근거는 hardware/gear/layout.md 3.7절과 docs/decisions/0002 (구동부 360, 벨트 HTD 5M, 축간 약 104).
% 원은 피치원이고 이는 그리지 않았다. (나)는 축 방향(가로)을 0.7배로 줄여 그렸다.
\begin{tikzpicture}[x=0.022cm,y=0.022cm,
    dim/.style={<->,cgray,line width=0.4pt},
    dl/.style={font=\tiny,text=black,inner sep=1pt},
    lab/.style={font=\tiny,align=left,inner sep=1pt},
    pc/.style={draw=cblue,fill=cbluebg,fill opacity=0.5},
    one/.style={draw=corange,line width=0.8pt}]
  % ================= (가) 축 방향에서 본 모습 =================
  \node[font=\footnotesize\bfseries,anchor=west] at (-235,268) {(가) 축 방향에서 본 피치원};
  % 다른 층의 큰 원판 (점선)
  \draw[draw=cgreen,dashed] (0,39) circle (80);
  \draw[one,dashed] (0,0) circle (81);
  % 피치원
  \draw[pc] (0,-120) circle (39);
  \draw[pc] (0,0) circle (26);
  \draw[pc] (0,39) circle (22.5);
  \draw[one] (0,39) circle (13);
  \draw[pc] (0,81) circle (19.5);
  \draw[pc] (0,120) circle (19.5);
  \draw[draw=cgray,dashed] (0,120) circle (57.3);
  % 모터와 벨트 (HTD 5M 18:72, 축간 약 104)
  \draw[cgray,dashed] (0,120) -- (-73.5,193.5);
  \draw[draw=cgray,fill=clight] (-73.5,193.5) circle (14.3);
  \node[lab,anchor=east] at (-92,205) {모터 풀리 18잇수\\(크랭크 끝)};
  \node[dl,anchor=east] at (-42,150) {104};
  % 축
  \foreach \y in {-120,0,39,81,120} { \fill[cgray] (0,\y) circle (3); }
  % 축 이름 (왼쪽)
  \foreach \y in {-120,0,39,81,120} { \draw[cgray,line width=0.3pt] (-3,\y) -- (-100,\y); }
  \node[lab,anchor=east] at (-102,120) {F: 26잇수 (모듈 1.5)\\풀리 72잇수 (5M)};
  \node[lab,anchor=east] at (-102,81) {I: 아이들 26잇수};
  \node[lab,anchor=east] at (-102,39) {C: 30잇수 (모듈 1.5)\\한 이빨 (모듈 2)};
  \node[lab,anchor=east] at (-102,0) {M: 26잇수 (모듈 2)\\한 이빨 원판 (모듈 3)};
  \node[lab,anchor=east] at (-102,-120) {S: 26잇수 (모듈 3)};
  \node[lab,anchor=north,text=cgreen] at (-55,-48) {슬롯 원판 \O160};
  % 치수 (오른쪽)
  \foreach \y in {-120,0,39,81,120} { \draw[cgray,line width=0.3pt] (3,\y) -- (150,\y); }
  \draw[dim] (105,81) -- (105,120) node[midway,dl,fill=white] {39};
  \draw[dim] (105,39) -- (105,81) node[midway,dl,fill=white] {42};
  \draw[dim] (105,0) -- (105,39) node[midway,dl,fill=white] {39};
  \draw[dim] (105,-120) -- (105,0) node[midway,dl,fill=white] {120};
  \draw[dim] (140,0) -- (140,120) node[midway,dl,fill=white] {120};
  % ================= (나) 축을 따라 본 배치 =================
  % 칸: 뱅크 쪽 0-90, 옆판 A 90-98, 가운데 98-248, 옆판 B 248-256, 바깥 256-346
  \begin{scope}[shift={(215,0)},xscale=0.7]
    \node[font=\footnotesize\bfseries,anchor=west] at (-25,268) {(나) 축을 따라 본 칸 배치};
    % 옆판
    \draw[draw=cgray,fill=clight] (90,-175) rectangle (98,240);
    \draw[draw=cgray,fill=clight] (248,-175) rectangle (256,240);
    % 축
    \draw[shaft] (-20,120) -- (300,120);
    \draw[shaft] (104,81) -- (130,81);
    \draw[shaft] (98,39) -- (336,39);
    \draw[shaft] (-20,0) -- (252,0);
    \draw[shaft] (-20,-120) -- (252,-120);
    % 라인 축 커플링 (뱅크 쪽 칸)
    \foreach \y in {120,0,-120} { \draw[draw=cgray,fill=cgray] (4,\y-14) rectangle (44,\y+14); }
    % 자리올림 2 (뱅크 쪽 칸)
    \draw[one,fill=corangebg] (56,-81) rectangle (80,81);
    \draw[pc,fill opacity=1] (56,-159) rectangle (80,-81);
    % 26:30 기어열 (가운데 칸)
    \draw[pc,fill opacity=1] (110,100.5) rectangle (124,139.5);
    \draw[pc,fill opacity=1] (110,61.5) rectangle (124,100.5);
    \draw[pc,fill opacity=1] (110,16.5) rectangle (124,61.5);
    % 자리올림 1 (가운데 칸)
    \draw[one,fill=corangebg] (150,26) rectangle (176,52);
    \draw[pc,fill opacity=1] (150,-26) rectangle (176,26);
    % 모터 (옆판 B에 고정, 몸통은 가운데 칸 위쪽 뒤)
    \draw[draw=cgray,fill=white] (101,150.5) rectangle (248,236.5);
    \draw[cgray,line width=1.2pt] (256,193.5) -- (336,193.5);
    \draw[draw=cgray,fill=clight] (264,179.2) rectangle (286,207.8);
    % 풀리와 슬롯 원판 (바깥 칸)
    \draw[draw=cgray,fill=white] (264,62.7) rectangle (286,177.3);
    \draw[draw=cgreen,fill=cgreenbg] (322,-41) rectangle (326,119);
    \draw[draw=cgray,fill=white] (314,96) rectangle (334,112);
    % 이름
    \node[lab,anchor=north,text=corange] at (68,-162) {자리올림 2};
    \node[lab,anchor=north,text=corange] at (163,-30) {자리올림 1};
    \node[lab,anchor=north] at (117,14) {26:30};
    \node[lab,anchor=south] at (24,136) {커플링};
    \node[lab] at (174,193.5) {모터};
    \node[lab,anchor=north] at (275,60) {풀리};
    \node[lab,anchor=north,text=cgreen] at (324,-45) {슬롯 원판};
    \node[lab,anchor=west] at (350,104) {센서 3};
    \node[lab,anchor=west] at (350,193.5) {크랭크};
    \node[lab,anchor=east] at (-22,120) {F};
    \node[lab,anchor=east] at (-22,0) {M};
    \node[lab,anchor=east] at (-22,-120) {S};
    \node[lab,anchor=south] at (30,215) {$\leftarrow$ 열 모듈 쪽};
    % 칸 치수
    \foreach \x in {0,90,98,248,256,346} { \draw[cgray,line width=0.3pt] (\x,-180) -- (\x,-205); }
    \draw[dim] (0,-198) -- (90,-198) node[midway,below=1pt,dl] {90};
    \draw[dim] (98,-198) -- (248,-198) node[midway,below=1pt,dl] {150};
    \draw[dim] (256,-198) -- (346,-198) node[midway,below=1pt,dl] {90};
    \draw[cgray,line width=0.3pt] (0,-205) -- (0,-245);
    \draw[cgray,line width=0.3pt] (346,-205) -- (346,-245);
    \draw[dim] (0,-238) -- (346,-238) node[midway,dl,fill=white] {346 (모듈 폭 360)};
  \end{scope}
\end{tikzpicture}
